\section{Motivation}


Due to the rapid development and sales of drones, especially following the war in 
Ukraine, there is currently a significant need for robust drone detection systems 
(\cite{macedo2025drone}). The biggest threats associated with drones include their 
capability to carry small explosives, as well as chemical and biological weapons 
(\cite{seidaliyeva2024advances}). There is also a private security aspect, as drones 
can be used for espionage against both critical infrastructure and private persons 
(\cite{seidaliyeva2024advances}). Drones could also be used to smuggle illegal items 
and weapons across borders more discreetly (\cite{seidaliyeva2024advances}), and are 
becoming increasingly cheaper and accessible (\cite{seidaliyeva2024advances}). The 
scale of the problem is reflected in incident data: over 150 drone incidents were 
reported globally in 2023 alone (\cite{seidaliyeva2024advances}), with 38\% related 
to smuggling across borders, 28\% involving delivery of contraband into prisons, and 
18\% directly affecting airports through near-collisions or runway closures 
(\cite{seidaliyeva2024advances}). Notable incidents include the temporary closure of 
Gatwick Airport in May 2023 and a near-miss between a commercial aircraft and a drone 
at Stansted Airport in December 2022 (\cite{seidaliyeva2024advances}), illustrating 
the direct threat to civil aviation safety. These issues highlight the need for robust 
drone detection systems for both personal and societal security 
(\cite{seidaliyeva2024advances, macedo2025drone}).

 \vspace{1em}

The four main methods of detecting drones are radar sensors, Radio Frequency (RF) 
sensors, audio sensors, and camera sensors (\cite{seidaliyeva2024advances}). While 
radar systems offer the longest detection range of over 10 km (\cite{macedo2025drone}) 
and RF-based systems achieve classification accuracies approaching 99\% 
(\cite{macedo2025drone}), both require specialist hardware and come with significantly 
higher costs (\cite{macedo2025drone}). Camera-based systems can achieve precision 
above 90\% at speeds up to 170 FPS (\cite{macedo2025drone}), and acoustic sensors 
can be built with consumer-grade microphones at a cost of only a few hundred dollars 
(\cite{macedo2025drone}). This cost differential motivates the approach taken in this 
project, which combines acoustic and visual sensing to achieve reliable detection 
without specialist hardware (\cite{macedo2025drone, seidaliyeva2024advances}).

The aim of this project is to develop a low-cost drone detection system using only 
3D-printed and off-the-shelf components, making robust drone detection accessible 
to private individuals and organisations that cannot justify the cost of commercial 
solutions.

\section{Project description}

\subsection{Socio-economic aspects}


The total hardware cost of the system amounts to approximately NOK~8,503 
(Table~\ref{tab:equipment}), making it accessible to private individuals, 
small businesses, and municipalities that cannot justify the cost of commercial 
drone detection solutions. Commercial EO/IR camera systems range from approximately 
€9,000 to €72,000 (NOK~97,000--778,000), and radar-based systems from €27,000 to 
over €450,000 (NOK~292,000--4,860,000) (\cite{airsight2025buyersguide}), placing them out of 
reach for most non-military users. Acoustic detection systems, by contrast, can be 
implemented with consumer-grade microphones for as little as a few hundred dollars 
(\cite{macedo2025drone}), and camera-based systems can similarly be built at low cost 
using standard processing boards (\cite{macedo2025drone}). The proposed system uses 
exclusively off-the-shelf components and open-source software, with no recurring 
licensing costs, and the 3D-printed mechanical components further reduce cost and 
allow local fabrication without specialist manufacturing.

Beyond security applications, accessible drone detection carries broader 
socioeconomic benefits (\cite{seidaliyeva2024advances}). Critical infrastructure 
such as power stations, water treatment facilities, and communication towers is 
increasingly vulnerable to drone-based surveillance and sabotage 
(\cite{seidaliyeva2024advances}), yet many operators lack the budget for professional 
detection systems (\cite{macedo2025drone}). A low-cost solution based on off-the-shelf 
components lowers the barrier to entry significantly, enabling wider adoption across 
both public and private sectors (\cite{macedo2025drone}). For private individuals, 
the system addresses growing concerns around privacy violations and property 
surveillance by consumer drones (\cite{seidaliyeva2024advances}), which existing 
legislation alone has proven insufficient to prevent (\cite{seidaliyeva2024advances}). 
Furthermore, by relying entirely on 3D-printed parts and standard electronic 
components, the system could be locally produced and maintained without specialist 
knowledge, making it suitable for deployment in remote areas or developing regions 
where commercial alternatives are unavailable.

\subsection{Existing solutions}

Existing drone detection solutions can be grouped into four main categories: radar, 
Radio Frequency (RF), acoustic, and camera-based systems 
(\cite{seidaliyeva2024advances, macedo2025drone}). Radar systems provide the longest 
detection range, exceeding 10 km in some implementations (\cite{macedo2025drone}), 
but suffer from high false target rates exceeding 50\% in real-world airport tests 
(\cite{macedo2025drone}), and struggle to distinguish drones from birds 
(\cite{macedo2025drone, seidaliyeva2024advances}). RF-based systems achieve 
classification accuracies approaching 99\% for known drone protocols 
(\cite{macedo2025drone}), but cannot detect autonomous drones that do not emit 
control signals (\cite{seidaliyeva2024advances}), and require dense sensor networks 
of more than 20 nodes for reliable coverage of large areas (\cite{macedo2025drone}). 
Acoustic systems are low-cost and require no line of sight 
(\cite{seidaliyeva2024advances}), but their effective detection range is typically 
limited to 160--290 metres depending on background noise conditions 
(\cite{macedo2025drone}). Camera-based systems using deep learning --- particularly 
the YOLO family of models --- have demonstrated strong real-time detection 
performance, with precision exceeding 90\% at speeds up to 170 FPS 
(\cite{macedo2025drone}), though they are susceptible to poor lighting conditions 
and require a clear line of sight (\cite{seidaliyeva2024advances, macedo2025drone}). 
Multimodal approaches combining several of these modalities show the greatest 
robustness, with positioning errors below 1.5\% of detection range in some 
implementations (\cite{macedo2025drone}), motivating the acoustic-visual fusion 
approach taken in this project.

\subsection{Problem statement}
The overall goal of this project is to develop a low-cost drone detection 
system using exclusively off-the-shelf and 3D-printed components, combining 
acoustic and visual sensing modalities. The camera and computer vision 
subsystem addresses the visual detection component of this system, with 
the core challenge being to achieve reliable real-time drone detection 
and tracking on resource-constrained embedded hardware, without relying 
on expensive dedicated computing infrastructure.

The camera subsystem was designed to fulfil four interconnected objectives. 
First, to train and deploy a YOLOv8n object detection model capable of 
reliably identifying drones across a variety of drone types, backgrounds, 
and distances. Second, to deploy this model on a Raspberry Pi 5 with a 
Hailo AI HAT+ accelerator such that inference could run in real time at 
a frame rate suitable for tracking a moving drone. Third, to publish the 
pixel coordinates of detected drones over MQTT to the Arduino-based motor 
controller, enabling the pan-tilt camera platform to orient toward and 
follow the detected target. Fourth, to stream the live camera feed 
including detection overlays to the web-based GUI, providing the 
operator with a real-time visual of the detection area.

The intended operational flow of the full system places the camera 
subsystem as the second stage of a two-stage detection pipeline. The 
acoustic microphone array first estimates the direction of arrival of 
drone sound, orienting the camera platform toward the target area. 
Once the camera is pointing in approximately the right direction, 
the YOLO detection takes over for visual confirmation and precise 
localisation, with the detected pixel coordinates driving continuous 
motor adjustments to keep the drone within the camera frame. This 
coarse-to-fine strategy means the camera subsystem does not need 
to scan the full 360° field of view independently, relying instead 
on the acoustic subsystem to narrow the search space first.

A key design constraint throughout was cost — the entire camera 
subsystem, including the Raspberry Pi 5, Hailo AI HAT+, and camera 
module, was required to remain within a budget accessible to private 
individuals and small organisations. Performance targets were defined 
by what was achievable within this constraint rather than by 
predetermined specifications, with the goal being to demonstrate 
that real-time drone detection and tracking is feasible on low-cost 
embedded hardware.

The real-time inference and coordinate publishing were successfully 
demonstrated in a live laboratory test. However, the full closed-loop 
tracking — where motor movement is driven continuously by YOLO 
detections following the initial acoustic orientation — was not 
completed within the project timeframe due to issues with the motor 
control hardware. The video streaming to the GUI was successfully 
implemented.